% !TEX program = pdflatex
\documentclass[12pt,a4paper]{article}

\usepackage[T1]{fontenc}
\usepackage[margin=1in]{geometry}
\usepackage{xcolor}
\usepackage{hyperref}
\usepackage{setspace}
\usepackage{parskip}
\usepackage{enumitem}

\hypersetup{
  colorlinks=true,
  urlcolor=blue!70!black
}

\onehalfspacing
\pagestyle{empty}

\begin{document}

\begin{center}
{\Large\bfseries Personal Statement}\\[4pt]
{\large Olusegun Israel Odeyemi}\\[2pt]
{\normalsize UK Global Talent Visa --- Digital Technology (Exceptional Promise)}
\end{center}

\bigskip

I am a Security Engineer and Technical Founder based in Derby, applying under the Global Talent route in Digital Technology (Exceptional Promise). My work sits where cyber security, applied machine learning, decentralised finance, and regulatory technology meet, and it has already taken concrete form: a publicly accessible, Sepolia-backed compliance platform with measurable engineering depth and a credible commercial thesis.

\textbf{Why the United Kingdom, and why Derby}

Britain is, in my judgment, the most coherent place in the world to build the kind of system I am building. The FCA's clear posture on digital assets, the country's leadership in regulatory technology, and the depth of its fintech and cyber-security ecosystems combine to make it both a credible market and a credible research environment. Derby, where I already live, gives me proximity to the Midlands' growing tech corridor and to Nottingham and Birmingham, while keeping me close enough to London for industry engagement. I intend to remain in the United Kingdom long term, scaling AMTTP from public testnet demonstration to regulated commercial deployment, and continuing to contribute to UK research output in this space.

\textbf{What I have built: AMTTP}

AMTTP, the Anti-Money Laundering Transaction Trust Protocol, is a blockchain-native compliance system designed to make risk decisions \emph{before} transaction finality. Most established digital-asset compliance tooling operates after settlement, when funds have already moved and risk has crystallised. AMTTP closes that gap by combining machine-learning risk scoring, graph-based relationship analysis, sanctions screening, and policy evaluation inside the Ethereum confirmation window, so a transaction can be approved, escrowed, reviewed, or blocked before on-chain settlement.

The platform is not a slide-deck architecture. It is a working, publicly accessible Sepolia-backed environment that I designed and built end to end, comprising:
\begin{itemize}[nosep]
  \item a cross-platform Flutter consumer application supporting wallet interaction, pre-send Trust Check scoring, disputes, and privacy-proof management;
  \item a Next.js 14 enterprise ``War Room'' dashboard offering real-time alerting, decision explainability, and compliance reporting;
  \item a FastAPI compliance orchestrator that fans out to nine specialist microservices---ML Risk, Graph, Policy, Sanctions, Monitoring, GeoRisk, Integrity, Explainability, and ZK-NAF verification---behind a TypeScript Express oracle gateway;
  \item a blockchain layer of 38 AMTTP-authored Solidity files, deployed through UUPS upgradeable proxies and covering dispute resolution, cross-chain risk propagation, and Groth16 zero-knowledge verification; and
  \item TypeScript and Python SDK layers that I also authored.
\end{itemize}

The work has been carried out independently, without institutional funding, formal academic research supervision, or venture backing.

\textbf{Applied machine learning}

The machine-learning component is operational rather than decorative. I trained and evaluated knowledge-distilled student models on 625{,}168 Ethereum address samples, including 372 confirmed fraud cases---a class imbalance that mirrors real fraud-detection conditions. The production CPU-optimised XGBoost student records ROC-AUC 0.9898; the LightGBM student reaches ROC-AUC 0.9999997 with F1 0.9906; and the ensemble meta-learner reports ROC-AUC 0.9999998 with PR-AUC 0.9997. Distillation was a deliberate engineering choice, not an academic one: it lets the system run on commodity CPUs rather than expensive GPU infrastructure, which is what makes the platform commercially deployable.

\textbf{Privacy, cryptography, and security-by-design}

I designed AMTTP to be privacy-preserving from the outset. Through zkNAF, my zero-knowledge compliance framework, the system supports sanctions non-membership verification, risk proofs, and KYC-related confirmations without unnecessary data exposure, using Groth16 SNARK verification over BN254 elliptic-curve pairings. This addresses a genuine commercial tension---meeting compliance obligations while minimising data collection---and aligns with FCA expectations and FATF Travel Rule principles. The same security-by-design discipline runs through the contract suite: Slither static analysis, Echidna property-based fuzzing at 50{,}000 iterations with 100-step sequences, and Hardhat integration tests covering deployment, transfer logic, policy enforcement, and cross-contract behaviour.

\textbf{Original research}

Alongside the platform, I have produced four research papers:
\begin{enumerate}[nosep]
  \item \emph{AMTTP Protocol} (TechRxiv), describing the system architecture;
  \item \emph{Blind-Spot Decomposition Theory (BSDT)} (Zenodo), an original mathematical framework for hidden model failure modes;
  \item \emph{Universal Detection Principle (UDP)} (Zenodo), a multi-operator extension of BSDT for anomaly detection; and
  \item \emph{The Geometry of System Collapse} (Zenodo), a phase-transition treatment of coupled-system instability validated against banking, Terra/Luna, and ERCOT grid data.
\end{enumerate}

The first two papers, which have had the longest exposure, have together attracted 365 views and 210 downloads as of 5 May 2026, achieved without institutional promotion. BSDT itself demonstrates classifier-blind-spot correction without retraining, with recall improvements of up to 84.4 percentage points and a flagship 99.94\% false-positive reduction on the Shuttle benchmark. I have kept BSDT as a separate research and optional analytics layer rather than embedding it in the default AML decision path, because the production student models are already strong and I did not want to introduce avoidable regulatory complexity into the core workflow. In April 2026, three of these papers were adopted as recommended readings in the postgraduate course CPE 604 by Dr Sunday Adeola Ajagbe, an independent researcher with no prior supervisory relationship to me---evidence that the work has begun to circulate beyond self-publication.

\textbf{Intellectual property and commercial intent}

On 4 March 2026 I filed UK patent application No.\ 1026066039, an eight-subsystem unified filing covering machine-learning risk scoring, blockchain compliance architecture, zero-knowledge identity workflows, cross-chain propagation, graph anomaly detection, adaptive friction, policy automation, and system-stability analysis. AMTTP is positioned as pre-settlement Compliance-as-a-Service infrastructure for regulated institutions and Virtual Asset Service Providers entering decentralised finance, and it already exposes consumer and enterprise interfaces, role-based access control, explainability tooling, and cross-chain capability through the public demonstration at \texttt{amttp.com}.

\textbf{Why this profile fits Exceptional Promise}

The relevant question for this route is not whether I am already a globally established figure, but whether I show clear potential to become a leading talent in UK digital technology. I believe the evidence supports that conclusion. I have built across Flutter, Next.js, Python microservices, Solidity, cryptography, and machine learning. I have translated original research into a working, demonstrable product. I have done so with the testing, auditing, and privacy posture that regulated infrastructure requires. And I have contributed to the UK ecosystem through publication, patent activity, and product-led innovation in a strategically important sector.

I am seeking endorsement so that I can continue this work in the United Kingdom---advancing AMTTP from public testnet demonstration toward regulated deployment, sustaining the associated research programme, and contributing to Britain's leadership in secure, intelligent, and commercially relevant digital infrastructure.

\end{document}
